\documentclass[11pt]{article}

% --------------------------------------------------
% Packages
% --------------------------------------------------
\usepackage[T1]{fontenc}
\usepackage{lmodern}
\usepackage{geometry}
\usepackage{microtype}
\usepackage{setspace}
\usepackage{amsmath,amssymb,amsthm}
\usepackage{csquotes}
\usepackage{hyperref}

\geometry{margin=1in}
\setstretch{1.15}

% --------------------------------------------------
% Theorem environments
% --------------------------------------------------
\newtheorem{definition}{Definition}
\newtheorem{proposition}{Proposition}
\newtheorem{remark}{Remark}

% --------------------------------------------------
% Title
% --------------------------------------------------
\title{The Limits of Interpretability}
\author{Flyxion}
\date{\today}

\begin{document}
\maketitle

% ==================================================
\begin{abstract}
Mechanistic interpretability has increasingly adopted causal language to move beyond correlational inspection of neural networks toward claims about the algorithms they implement. This essay presents a unified analysis of this causal paradigm. We examine the role of interventions and counterfactuals, the criteria for mechanistic responsibility, the construction and validation of causal abstractions, and the interpretation of failure modes. We argue that the resulting framework is methodologically powerful yet deliberately limited: it privileges present-state interventions over irreducible history, treats irreversibility as informational rather than ontological, admits plural and domain-restricted abstractions, and defines understanding as successful reconstruction under intervention rather than exhaustive description. The essay concludes by clarifying what this paradigm can explain, what it necessarily leaves open, and why these limits are not defects but structural features of causal explanation itself.
\end{abstract}

% ==================================================
\section{Introduction}
% ==================================================

The rapid progress of large neural networks has intensified demands for explanations that go beyond performance metrics and surface-level inspection. In response, a causal paradigm of mechanistic interpretability has emerged, aiming not merely to observe internal activations, but to determine which internal components are responsible for which aspects of behavior. This paradigm seeks to answer a specific question: under what conditions can a neural system be said to implement a particular algorithm?

The distinguishing feature of this approach is its reliance on interventions and counterfactual reasoning. Rather than inferring mechanism from correlation, the framework evaluates hypotheses by actively modifying internal variables and observing the consequences. The goal is not descriptive transparency, but causal accountability.

This essay develops a systematic account of this paradigm. We begin by clarifying its conception of mechanisms and explanatory units. We then analyze how causality, abstraction, and failure interact within the framework, and finally assess the scope and limits of the resulting notion of understanding.

% ==================================================
\section{Mechanisms and the Status of Internal States}
% ==================================================

A defining commitment of causal mechanistic interpretability is the treatment of internal model components as causal variables. These variables are typically instantiated as instantaneous values of activations at specific layers or units. Interventions act by setting such values to quantities they would not naturally assume and measuring the resulting changes in output.

\begin{definition}[Intervention]
An intervention is an operation that sets an internal variable of a model to a specified value independently of its usual causal antecedents.
\end{definition}

This definition privileges the present state of the system as the locus of causal testing. The explanatory unit is the value being set, not the historical process by which it was produced. However, this does not entail that history is irrelevant. Distinct causal lineages can be distinguished if and only if they manifest as differences in downstream behavior under counterfactual substitution.

Thus, the framework adopts a conditional stance toward history: historical differences matter precisely when they cash out into present causal effects. There is no commitment to irreducible historical identity beyond what is functionally testable.

% ==================================================
\section{Causality, Counterfactuals, and Responsibility}
% ==================================================

Causal responsibility within this paradigm is established through mediation rather than mere association. A component is said to bear responsibility for a behavior if manipulating it alters that behavior in a systematic and predictable way.

One prominent technique involves degrading performance by corrupting a representation and then selectively restoring candidate mediators to their original values. If restoring a component suffices to recover the behavior, that component is identified as causally relevant.

Importantly, responsibility is rarely localized. Instead, it is distributed across structured pathways involving multiple components whose roles are defined relationally.

\begin{proposition}
Mechanistic sufficiency in neural systems is typically relational rather than atomic: components mediate behavior only within structured causal contexts.
\end{proposition}

Moreover, causality in this setting is often probabilistic. Gradual degradation under noise and partial restoration under intervention indicate that mechanisms can be real without being deterministically necessary or sufficient in isolation.

% ==================================================
\section{Interchange Interventions and Semantic Substitution}
% ==================================================

A central methodological innovation of the causal paradigm is the interchange intervention, in which an internal value from one input context is substituted into another. This operation tests whether a component plays a stable causal role across contexts.

There is no a priori guarantee that such substitutions are semantically coherent. Rather, coherence is determined empirically: if the substitution produces the expected change in behavior, the hypothesis is supported. If it produces novel or incoherent behavior, the proposed causal model must be revised.

This renders semantic compatibility an operational notion. Meaning is not assumed; it is inferred from successful transport across contexts. However, this also implies that failures of substitution diagnose limitations of the hypothesized abstraction rather than revealing deeper semantic facts on their own.

% ==================================================
\section{Causal Abstraction and Algorithmic Explanation}
% ==================================================

The primary explanatory ambition of the framework is causal abstraction. A causal abstraction relates a low-level system to a higher-level algorithm by establishing a systematic correspondence between interventions at the two levels.

\begin{definition}[Causal Abstraction]
A causal abstraction is a mapping between two causal models such that corresponding interventions produce equivalent effects on their respective outputs.
\end{definition}

Crucially, multiple abstractions can be valid simultaneously. Different high-level descriptions may capture different regularities of the same system, provided they preserve interventional equivalence within their domains. There is no unique privileged abstraction, nor an expectation that a single algorithm governs all behaviors.

An abstraction is considered excessive not when it is false, but when it ceases to provide explanatory compression or interpretive clarity.

% ==================================================
\section{Irreversibility and Information Loss}
% ==================================================

Irreversibility enters the framework not as a fundamental principle, but as a consequence of information destruction. Interventions such as ablation or noise injection can degrade representations in ways that cannot be fully undone by subsequent manipulations.

However, this irreversibility is epistemic rather than ontological. It reflects the loss of recoverable signal, not an intrinsic asymmetry of causation. As a result, collapse is interpreted as a boundary of explanation rather than a basic feature of mechanism.

This distinction allows the framework to remain agnostic about deeper metaphysical claims while retaining practical criteria for explanatory adequacy.

% ==================================================
\section{Failure Modes and Algorithmic Mixtures}
% ==================================================

Empirical application of causal interpretability methods often reveals that abstractions hold only under restricted conditions. Performance may degrade gracefully in some regimes and collapse entirely in others.

Such failures are not merely negative results. They provide evidence that the system may implement a mixture of heuristics rather than a single coherent algorithm. In these cases, explanation may be necessarily patchwork, with different causal stories applying in different contexts.

The framework thus accommodates partial validity and domain restriction as expected outcomes rather than anomalies.

% ==================================================
\section{Understanding Without Finality}
% ==================================================

Within this paradigm, understanding is achieved when a system can be reconstructed as a simpler causal process that predicts behavior under intervention. Control and prediction are valuable, but they are subordinate to the identification of internal causal structure.

At the same time, successful abstraction does not imply completeness. Any abstraction is a zoomed-out view that omits detail. Passing all current tests does not guarantee that no further structure remains undiscovered.

Understanding is therefore provisional, layered, and revisable by design.

% ==================================================
\section{Conclusion}
% ==================================================

The causal paradigm of mechanistic interpretability offers a rigorous methodology for transforming neural networks from opaque function approximators into systems that admit algorithmic explanation. By grounding claims in interventions and counterfactuals, it replaces metaphor with testable structure.

At the same time, its limits are intrinsic. It privileges present-state causality over irreducible history, treats irreversibility as informational loss, admits plural and context-dependent abstractions, and defines understanding as successful causal reconstruction rather than exhaustive description.

These limits should not be read as shortcomings. They reflect a disciplined refusal to claim more than causal evidence can support. In doing so, the paradigm establishes a clear boundary between what can be reverse-engineered, what can be controlled, and what must remain open to further inquiry.

% ==================================================
\section{Formal Structure of Causal Mechanistic Explanation}
% ==================================================

This section formalizes the causal paradigm described above by specifying its basic objects, mappings, and validation criteria. The aim is not to introduce new metaphysical commitments, but to make explicit the logical structure implicit in intervention-based mechanistic explanation.

% --------------------------------------------------
\subsection{Causal Models and Variables}
% --------------------------------------------------

Let a neural system be represented as a causal model
\[
\mathcal{M}_L = (V_L, E_L, F_L),
\]
where:
\begin{itemize}
\item $V_L$ is a finite set of low-level variables (e.g., internal activations at specified layers or components),
\item $E_L \subseteq V_L \times V_L$ is a directed acyclic graph encoding causal dependencies,
\item $F_L$ is a set of structural equations assigning each variable a value as a function of its parents.
\end{itemize}

An intervention $\mathrm{do}(X := x)$ replaces the structural equation for $X \in V_L$ with the constant assignment $x$, yielding a modified model $\mathcal{M}_L^{\mathrm{do}(X:=x)}$.

Outputs of the system are designated variables $Y \subseteq V_L$.

% --------------------------------------------------
\subsection{High-Level Causal Descriptions}
% --------------------------------------------------

A candidate algorithmic explanation is represented as a higher-level causal model
\[
\mathcal{M}_H = (V_H, E_H, F_H),
\]
whose variables correspond to abstract computational roles (e.g., carries, gates, retrieved facts, or control signals).

The high-level model is not assumed to be isomorphic to the low-level model. Instead, it is connected to it via a mapping.

\begin{definition}[Abstraction Mapping]
An abstraction mapping is a function
\[
\alpha : V_L \rightarrow V_H \cup \{\bot\},
\]
assigning each low-level variable either to a high-level variable or to a null symbol indicating irrelevance at the chosen level of abstraction.
\end{definition}

Variables mapped to $\bot$ are treated as internal detail that is ignored rather than denied.

% --------------------------------------------------
\subsection{Interventional Correspondence}
% --------------------------------------------------

The central criterion of causal abstraction is interventional correspondence.

\begin{definition}[Interventional Correspondence]
A low-level intervention $\mathrm{do}(X := x)$ corresponds to a high-level intervention $\mathrm{do}(Z := z)$ if
\[
\alpha(X) = Z
\quad \text{and} \quad
\alpha(x) = z,
\]
where $\alpha(x)$ denotes the induced value of the abstract variable.
\end{definition}

The abstraction is evaluated by comparing the effects of corresponding interventions on their respective outputs.

\begin{definition}[Causal Abstraction]
A high-level model $\mathcal{M}_H$ is a causal abstraction of $\mathcal{M}_L$ under $\alpha$ if, for all corresponding interventions,
\[
P(Y_L \mid \mathrm{do}(X := x)) \approx P(Y_H \mid \mathrm{do}(Z := z)),
\]
where $Y_L$ and $Y_H$ are the designated outputs of the low- and high-level models respectively.
\end{definition}

The approximation reflects the empirical and probabilistic nature of neural systems; exact equality is not required.

% --------------------------------------------------
\subsection{Mechanistic Responsibility}
% --------------------------------------------------

Causal responsibility is defined via mediation.

\begin{definition}[Causal Mediator]
A variable $M \in V_L$ is a mediator of the effect of input $I$ on output $Y$ if intervening on $M$ alters the conditional distribution
\[
P(Y \mid \mathrm{do}(I := i)).
\]
\end{definition}

Restoration experiments operationalize this notion: if corrupting the input degrades performance and restoring $M$ recovers it, $M$ is identified as causally responsible relative to the tested behavior.

Responsibility is therefore relational and path-dependent rather than intrinsic to isolated variables.

% --------------------------------------------------
\subsection{Granularity and Plurality of Abstractions}
% --------------------------------------------------

Nothing in the formalism enforces uniqueness of abstraction.

\begin{proposition}
Distinct abstraction mappings $\alpha_1, \alpha_2$ may each induce valid causal abstractions of the same low-level model, provided they satisfy interventional correspondence within their respective domains.
\end{proposition}

This formalizes the plurality of explanations: different abstractions compress different regularities while discarding different details.

% --------------------------------------------------
\subsection{Failure and Domain Restriction}
% --------------------------------------------------

Let $\mathcal{D}$ be a distribution over inputs.

\begin{definition}[Domain-Restricted Validity]
A causal abstraction is valid on $\mathcal{D}$ if interventional correspondence holds for all tested interventions drawn from $\mathcal{D}$.
\end{definition}

Failure outside $\mathcal{D}$ does not falsify the abstraction; it restricts its scope. Empirically observed collapse may therefore indicate either representational destruction or the presence of multiple internal strategies not captured by a single high-level model.

% --------------------------------------------------
\subsection{Limits of Formal Reconstruction}
% --------------------------------------------------

The formalism deliberately avoids treating irreversibility as primitive. Information-destroying interventions may render certain causal pathways unrecoverable, but this reflects loss of signal rather than violation of the causal structure itself.

Moreover, satisfaction of the abstraction criteria does not imply completeness.

\begin{remark}
A system may admit a valid causal abstraction at one level while containing further internal structure at finer scales that remain unexplained.
\end{remark}

The formal framework therefore characterizes mechanistic understanding as the construction of reliable interventional models, not as an exhaustive account of all internal organization.

% ==================================================
\section{Categorical and Operational Semantics Perspectives}
% ==================================================

The causal framework developed above can be expressed equivalently in categorical and operational-semantic terms. Doing so clarifies which aspects of the account are structural, which are empirical, and which are deliberately left uninterpreted. This section provides two complementary reformulations.

% --------------------------------------------------
\subsection{Causal Models as Structured Morphisms}
% --------------------------------------------------

We begin by treating causal models as objects in a category.

Let $\mathbf{C}$ be a category whose objects are causal models
\[
\mathcal{M} = (V, E, F),
\]
and whose morphisms preserve causal structure up to abstraction. Intuitively, a morphism represents a structure-preserving map between explanatory descriptions.

\begin{definition}[Causal Morphism]
A causal morphism
\[
\phi : \mathcal{M}_L \rightarrow \mathcal{M}_H
\]
is a mapping that sends variables in $V_L$ to variables in $V_H \cup \{\bot\}$ and preserves directed dependencies among variables that are not mapped to $\bot$.
\end{definition}

Such morphisms need not be injective or surjective. Many low-level variables may collapse into a single high-level variable, and many low-level details may be discarded entirely.

In this setting, a causal abstraction is not an isomorphism but a homomorphism: it preserves relevant structure while intentionally forgetting detail.

% --------------------------------------------------
\subsection{Interventions as Endomorphisms}
% --------------------------------------------------

Interventions can be represented categorically as endomorphisms acting on causal models.

\begin{definition}[Intervention Endomorphism]
An intervention $\mathrm{do}(X := x)$ induces an endomorphism
\[
\iota_{X := x} : \mathcal{M} \rightarrow \mathcal{M}
\]
that replaces the structural equation for $X$ with a constant assignment.
\end{definition}

Interventional correspondence between abstraction levels can then be expressed as commutativity of diagrams.

\begin{definition}[Interventional Commutativity]
A diagram
\[
\begin{array}{ccc}
\mathcal{M}_L & \xrightarrow{\iota_{X := x}} & \mathcal{M}_L \\
\downarrow \phi &  & \downarrow \phi \\
\mathcal{M}_H & \xrightarrow{\iota_{Z := z}} & \mathcal{M}_H
\end{array}
\]
commutes if the effects of the low-level and high-level interventions correspond under $\phi$.
\end{definition}

Causal abstraction is precisely the requirement that such diagrams commute approximately for the class of tested interventions.

% --------------------------------------------------
\subsection{Plurality as Non-Uniqueness of Factorizations}
% --------------------------------------------------

From this perspective, the existence of multiple valid abstractions follows immediately.

\begin{proposition}
If a causal model admits multiple distinct morphisms
\[
\phi_1, \phi_2 : \mathcal{M}_L \rightarrow \mathcal{M}_H^{(1)}, \mathcal{M}_H^{(2)}
\]
that satisfy interventional commutativity, then the system admits multiple valid abstractions.
\end{proposition}

No universal property selects a unique factorization. Explanatory adequacy is therefore indexed to the chosen morphism and its domain of validity.

% --------------------------------------------------
\subsection{Operational Semantics of Mechanistic Explanation}
% --------------------------------------------------

We now give an operational reading of the same framework.

Let $\Sigma$ denote the space of internal states of a model, and let $\mathcal{O}$ denote the space of observable outputs. A model execution is a transition
\[
\sigma \xrightarrow{\,\mathsf{step}\,} \sigma'
\quad \text{with output} \quad o \in \mathcal{O}.
\]

An intervention is an operation on states:
\[
\mathsf{set}_{X := x} : \Sigma \rightarrow \Sigma,
\]
which overwrites the value of a designated component.

\begin{definition}[Operational Counterfactual]
Given an execution trace $\sigma \rightarrow o$, a counterfactual execution is obtained by composing
\[
\sigma' = \mathsf{set}_{X := x}(\sigma)
\]
and evaluating the resulting output $o'$.
\end{definition}

Mechanistic claims are validated by comparing the operational behavior of original and counterfactual executions.

% --------------------------------------------------
\subsection{Operational Meaning of Abstraction}
% --------------------------------------------------

A high-level description induces a quotient of the state space.

\begin{definition}[Abstract State Projection]
An abstraction is a projection
\[
\pi : \Sigma \rightarrow \Sigma_H
\]
that maps concrete states to abstract states, discarding irrelevant detail.
\end{definition}

The abstraction is valid if projected executions simulate abstract executions:

\begin{definition}[Operational Soundness]
An abstraction $\pi$ is sound if for all tested interventions,
\[
\pi(\mathsf{step}(\mathsf{set}_{X := x}(\sigma))) \approx
\mathsf{step}_H(\mathsf{set}_{Z := z}(\pi(\sigma))).
\]
\end{definition}

This expresses causal abstraction as a simulation relation rather than an identity.

% --------------------------------------------------
\subsection{Failure, Collapse, and Non-Simulation}
% --------------------------------------------------

Operationally, failure corresponds to breakdown of simulation.

If an intervention yields a state for which no coherent abstract transition exists, the abstraction ceases to be sound in that regime. Collapse thus indicates not an error of execution, but a mismatch between the chosen abstraction and the system’s operational dynamics.

Importantly, nothing in this semantics requires global soundness. Simulation may hold locally and fail elsewhere.

% --------------------------------------------------
\subsection{What the Semantics Do Not Fix}
% --------------------------------------------------

Both the categorical and operational formulations deliberately leave several questions open:
\begin{itemize}
\item which abstraction is preferable,
\item whether abstractions should compose,
\item whether history should be encoded explicitly in states,
\item and whether irreversibility should be treated as primitive.
\end{itemize}

These are not oversights but boundaries. The formalism specifies how mechanistic claims are tested and related, not which claims must ultimately be made.

\begin{remark}
From both perspectives, mechanistic understanding is the construction of structure-preserving maps that make counterfactual behavior predictable, not the recovery of a unique or final description.
\end{remark}

\appendix
% ==================================================
\section{Appendix: Operational Realization of Spherepop Operators}
% ==================================================

This appendix presents an operational realization of Spherepop operators as an
event-historical calculus. The purpose is not to privilege any particular
computational substrate, but to demonstrate that the Spherepop semantics
developed in the main text admit a direct, inspectable realization in which
events are authoritative and all present structure is derived.

The guiding constraint throughout is \emph{event-sourcing}: the semantic ground
truth of the system is a log of irreversible events, while any notion of
``current state'' is a projection obtained by interpreting (replaying) that log.
No mutable state is treated as primitive. Instead, state exists only as a
convenience cache whose validity is always subordinate to replay equivalence.

From a semantic perspective, a history is a morphism in a freely generated
temporal category, and replay is a functor from that category into a concrete
state space. Two histories are equivalent exactly insofar as this functor maps
them to the same derived invariants. Collapse is therefore not a mutation of
state but a transformation of the history itself: it alters which distinctions
remain available for counterfactual explanation while preserving the induced
present.

Operationally, the system is structured around three layers:

\begin{enumerate}
\item \textbf{Event layer.} A closed set of event constructors defines the only
irreversible operations. These events constitute the sole source of semantic
authority.
\item \textbf{Replay semantics.} A total interpreter maps histories to derived
structure. This interpreter is pure: its output depends only on the event
sequence, not on ambient context.
\item \textbf{World interface.} A convenience interface maintains both a history
and a derived cache, ensuring that every user-facing operation appends an event
and that the cache remains replay-consistent.
\end{enumerate}

This separation mirrors the distinction developed in the main text between
event-nodes and texture-nodes. Event-nodes appear explicitly in the log and act
on the base category of admissible futures. Texture-nodes, by contrast, are not
logged; they appear only as parameters to interpretation or event proposal, and
remain fiberwise so long as they preserve future support and reconstructibility.

From a categorical standpoint, replay defines a functor from the free category
generated by events into a category of derived states. Collapse corresponds to a
quotienting of morphisms that preserves functorial image while erasing internal
factorization. From an operational-semantics standpoint, replay is a fold over
events, and counterfactual reasoning corresponds to replay under modified
prefixes or substituted events.

The constructions below should therefore be read neither as an implementation
guide nor as a commitment to any particular runtime model, but as an existence
proof: the Spherepop calculus is not merely conceptual. It can be realized as a
small, explicit, replayable event algebra in which irreversibility, commitment,
and abstraction are expressed directly in structure rather than implicitly in
control flow.


\end{document}
